% 沿用本 sandbox 既有中文单栏模板，仅保留本论文所需包。
\usepackage[letterpaper,margin=1in]{geometry}
\usepackage{lmodern,fontspec}
\setmainfont{TeX Gyre Pagella}
\setCJKmainfont{FandolSong-Regular.otf}[BoldFont=FandolSong-Bold.otf,ItalicFont=FandolKai-Regular.otf,Script=Default]
\setCJKsansfont{FandolHei-Regular.otf}[BoldFont=FandolHei-Bold.otf,Script=Default]
\setCJKmonofont{FandolFang-Regular.otf}[Script=Default]
\usepackage{microtype,mathtools,amssymb,bm}
\usepackage{graphicx,booktabs,enumitem,caption,subfig,float,adjustbox}
\usepackage{placeins,needspace}
\usepackage{xcolor,titlesec,xurl}
\definecolor{paperblue}{HTML}{1F4E79}
\usepackage[numbers,square,compress]{natbib}
\usepackage{multirow}
\usepackage[colorlinks=true,linkcolor=paperblue,citecolor=paperblue,urlcolor=paperblue,bookmarksnumbered=true]{hyperref}
\special{pdf:minorversion 7}
% 论文特定的 PDF metadata 在填写正文时设置，不残留上一论文的标题/作者。
\renewcommand{\eqref}[1]{(\ref{#1})}
\renewcommand{\refname}{参考文献}
\renewcommand{\abstractname}{摘要}
\setlength{\parindent}{2em}
\setlength{\parskip}{0.35em}
\setlength{\emergencystretch}{2em}
\setlength{\tabcolsep}{1.8mm}
\renewcommand{\arraystretch}{1.15}
\setlist[enumerate]{leftmargin=2.2em,itemsep=0.2em}
\captionsetup{font=small,labelfont={bf,color=paperblue},labelsep=quad}
\titleformat{\section}{\Large\sffamily\bfseries\color{paperblue}}{\thesection}{0.75em}{}
\titleformat{\subsection}{\large\sffamily\bfseries}{\thesubsection}{0.65em}{}
\titlespacing*{\section}{0pt}{1.5em}{0.55em}
\titlespacing*{\subsection}{0pt}{1em}{0.35em}
\renewcommand{\topfraction}{0.9}
\renewcommand{\bottomfraction}{0.8}
\renewcommand{\textfraction}{0.08}
\renewcommand{\floatpagefraction}{0.75}
\setcounter{topnumber}{3}
\setcounter{totalnumber}{5}
\makeatletter\setlength{\@fptop}{0pt}\makeatother
